\documentclass[tcc,capa]{texifsap}

\usepackage[utf8]{inputenc} % acentuacao
\usepackage{graphicx} % para inserir figuras
\usepackage[T1]{fontenc}


\hypersetup{
    hidelinks, % Remove coloração e caixas
    unicode=true,   %Permite acentuação no bookmark
    linktoc=all %Habilita link no nome e página do sumário
}

\unidade{C{\^a}mpus Sapucaia do Sul}

%TÍTULO DE FORMAÇÃO: descomente o bloco correspondente ao seu curso e
%mantenha o outro comentado.
%
%--- Curso Superior de Tecnologia em Análise e Desenvolvimento de Sistemas:
\curso{An\'alise e Desenvolvimento de Sistemas}
\nomecurso{Superior de Tecnologia em An\'alise e Desenvolvimento de Sistemas}
%na capa, o nome do curso quebra a linha no ponto indicado pelo \\
\nomecursocapa{Superior de Tecnologia em\\An\'alise e Desenvolvimento de Sistemas}
\titulocurso{Tecn\'ologo(a) em An\'alise e Desenvolvimento de Sistemas}
%
%--- Curso Técnico em Desenvolvimento de Sistemas:
%\curso{Desenvolvimento de Sistemas}
%\nomecurso{T\'ecnico em Desenvolvimento de Sistemas}
%\titulocurso{T\'ecnico(a) em Desenvolvimento de Sistemas}

\title{Um Blabla Blablabla com Aplicações em Blablabla}

\author{ÚltimoSobrenome1}{Primeiro nome e sobrenome}
\author{ÚltimoSobrenome2}{Primeiro nome e sobrenome}                   %comente caso seja um trabalho em duplas
\advisor[Prof.~Dr.]{ÚltimoSobrenome}{Primeiro nome e sobrenome}
\coadvisor[Prof.~Dr.]{ÚltimoSobrenome}{Primeiro nome e sobrenome}      %comente caso tenhas co-orientador

\begin{document}

%\renewcommand{\advisorname}{Orientadora}           %descomente caso tenhas orientadora
%\renewcommand{\coadvisorname}{Coorientadora}      %descomente caso tenhas co-orientadora

%Palavras-chave em PT_BR
%REGRA (NBR 6028:2021): de 3 a 5 palavras-chave, grafadas com inicial
%minúscula (exceto nomes próprios), separadas entre si por ponto e vírgula
%e com a última finalizada por ponto. Ex.: \keyword{engenharia de software;}
\keyword{palavrachave-um;}
\keyword{palavrachave-dois;}
\keyword{palavrachave-tres;}
\keyword{palavrachave-quatro;}
\keyword{palavrachave-cinco.}

%Palavras-chave em EN_US
\keywordeng{keyword-one;}
\keywordeng{keyword-two;}
\keywordeng{keyword-three;}
\keywordeng{keyword-four;}
\keywordeng{keyword-five.}


%\renewcommand{\advisorname}{Orientadora}           %descomente caso tenhas orientadora
%\renewcommand{\coadvisorname}{Coorientadora}      %descomente caso tenhas coorientadora

\maketitle

\sloppy

%FICHA CATALOGRÁFICA (elemento obrigatório na VERSÃO FINAL):
%a página gerada abaixo é o verso da folha de rosto e fica reservada
%para a ficha catalográfica. A ficha NÃO é elaborada pelo estudante:
%solicite-a à Biblioteca do câmpus DEPOIS da banca e das correções
%finais (dados como o número de folhas mudam até lá). Ao receber o
%PDF da Biblioteca (ex.: ficha.pdf), salve-o nesta pasta e troque a
%linha \fichacatalografica por:
%\fichacatalograficapdf{ficha.pdf}
%Conforme a NBR 14724:2024, esta página não é contada nem numerada —
%o template já cuida disso automaticamente.
\fichacatalografica

%Composição da Banca Examinadora: descomente abaixo para a versão final do TCC
%PRAZOS (Regulamento do TCC, arts. 24 e 30): a cópia do TCC deve ser
%enviada aos membros da Banca Examinadora com, no mínimo, 15 dias de
%antecedência da defesa; a versão final corrigida deve ser entregue no
%prazo definido pela banca, de no máximo 10 dias após a apresentação
%oral.

\begin{aprovacao}{30 de novembro de 2025} %data da banca por extenso
\noindent Prof(a). Nome completo,\\
Orientador(a)\\
<Nome da instituição>\\[1cm]
\noindent Prof(a). Nome completo,\\
Avaliador(a)\\
<Nome da instituição>\\[1cm]
\noindent Prof(a). Nome completo,\\
Avaliador(a)\\
<Nome da instituição>\\[1cm]
\noindent Prof(a). Nome completo,\\
Avaliador(a)\\
<Nome da instituição>\\[1cm]
\end{aprovacao}

%Opcional
\begin{dedicatoria}
  Dedico\ldots 
\end{dedicatoria}

%Opcional
\begin{agradecimentos}
  Agradeço\ldots 
\end{agradecimentos}

%Opcional
\begin{epigrafe}
  Só sei que nada sei.\\
  {\sc --- Sócrates}
\end{epigrafe}

%Resumo em Portugues
%REGRA (NBR 6028:2021): parágrafo ÚNICO, de 150 a 500 palavras, sem
%recuo, sem citações, sem siglas não definidas e sem fórmulas/tabelas.
%Redigir com verbo na voz ativa e na terceira pessoa do singular,
%iniciando com o objetivo do trabalho e apresentando método, resultados
%e conclusões.
\begin{abstract}
Escreva aqui o resumo do trabalho, em um único parágrafo de 150 a 500
palavras, sem recuo e sem citações: comece pelo objetivo do trabalho e
apresente, na sequência, o método utilizado, os principais resultados e
a conclusão. O resumo é a porta de entrada do trabalho: muitos leitores
decidem continuar (ou não) a leitura a partir dele.
\end{abstract}

%Resumo em Inglês (Abstract)
%REGRA (NBR 6028:2021): tradução fiel do resumo em português, seguindo
%as mesmas regras (parágrafo único, 150 a 500 palavras).
\begin{englishabstract}{Titulo do Trabalho em Ingles}
Write here the English version of the resumo, as a faithful
translation: a single paragraph of 150 to 500 words, beginning with the
objective of the work, followed by the method, the main results and the
conclusion.
\end{englishabstract}

%Lista de Figuras
\listoffigures

%Lista de Quadros
\listofquadros

%Lista de Tabelas
\listoftables

%lista de abreviaturas e siglas
\begin{listofabbrv}{ABNT}%coloque aqui a maior sigla para ajustar a distância
        \item[ABNT] Associação Brasileira de Normas Técnicas
        \item[NUMA] Non-Uniform Memory Access
        \item[SIMD] Single Instruction Multiple Data
        \item[SMP] Symmetric Multi-Processor
        \item[SPMD] Single Program Multiple Data
\end{listofabbrv}

%Sumario
\tableofcontents

%==============================================================================
% PARTE TEXTUAL — organização recomendada do TCC para pesquisas na área
% de computação (alinhada ao guia exemplo-tcc.docx):
%   1 Introdução (com Objetivo Geral e Objetivos Específicos)
%   2 Fundamentação Teórica
%   3 Trabalhos Relacionados
%   4 Metodologia
%   [5 Desenvolvimento — apenas para trabalhos de desenvolvimento de um
%      sistema. Descomente o capítulo correspondente adiante]
%   5/6 Resultados
%   Conclusão, Referências, Apêndices e Anexos
%==============================================================================

%ATENÇÃO — PLÁGIO E INTELIGÊNCIA ARTIFICIAL (Regulamento do TCC, art. 34):
%plágio total ou parcial ANULA o TCC e reprova o acadêmico. O uso de
%ferramentas de inteligência artificial para a produção do trabalho é
%considerado plágio. Escreva o texto com suas próprias palavras e cite e
%referencie TODAS as fontes utilizadas, conforme as regras de citação
%indicadas na Fundamentação Teórica.

\chapter{Introdução}

%GUIA: na introdução, situe o leitor no contexto do tema pesquisado,
%oferecendo uma visão global do estudo. Esclareça as delimitações
%trabalhadas na abordagem do assunto, os objetivos e as justificativas
%que levaram à investigação e, em seguida, aponte as questões de
%pesquisa para as quais se buscará resposta. Ao final da introdução,
%pode-se descrever sucintamente a forma como o trabalho está organizado.

Escreva aqui a introdução do trabalho. Situe o leitor no contexto do
tema pesquisado, oferecendo uma visão global do estudo, e esclareça as
delimitações adotadas na abordagem do assunto. Apresente o problema
investigado em linhas gerais, preparando o leitor para as seções
seguintes. No último parágrafo, ou na seção opcional Estrutura do
texto, ao final deste capítulo, descreva sucintamente a organização
do trabalho, indicando o que será tratado em cada capítulo.

\section{Problema de pesquisa}

%GUIA: o problema é o ponto de partida de todo o trabalho — objetivos,
%metodologia e resultados existem para resolvê-lo ou compreendê-lo.

Caracterize aqui o problema a ser pesquisado: descreva a situação
atual, as dificuldades, limitações ou lacunas observadas e as
consequências de não enfrentá-las, apoiando-se em dados e em fontes
citadas. Delimite também o contexto em que o problema será estudado
(onde, quando, com quem): um problema bem caracterizado e bem
delimitado é o que torna a pesquisa viável dentro do prazo do TCC.

\subsection{Pergunta de pesquisa}

A pergunta de pesquisa é a síntese do problema em forma interrogativa.
Formule-a aqui em uma única frase clara e delimitada, por exemplo:
``em que medida a técnica X melhora o aspecto Y no contexto Z?''. Todos
os capítulos seguintes devem contribuir para respondê-la.

\subsection{Hipótese de pesquisa - Opcional}

Enuncie aqui a resposta provisória à pergunta de pesquisa, que será
confirmada ou refutada pelos resultados do trabalho. A hipótese deve
ser verificável: evite afirmações vagas, que não possam ser testadas.
Por ser um elemento opcional, remova esta subseção caso o seu trabalho
não utilize hipótese.

\section{Justificativa}

Explique aqui por que o trabalho merece ser realizado: a relevância do
problema, os benefícios esperados para o público-alvo e para a área de
computação e as razões que motivaram a escolha do tema. Uma boa
justificativa apoia-se em dados e em fontes citadas, e não apenas na
opinião do autor.



\section{Objetivos}

Nas seções abaixo estão descritos o objetivo geral e os objetivos
específicos desta pesquisa.

\subsection{Objetivo geral}

%GUIA: um único objetivo, iniciado por verbo no infinitivo (analisar,
%desenvolver, avaliar...), coerente com o título do trabalho.

Enuncie aqui, em uma única frase iniciada por verbo no infinitivo
(analisar, desenvolver, avaliar, comparar\ldots), o resultado principal
que o trabalho pretende alcançar, coerente com o título do trabalho e
com a pergunta de pesquisa.

\subsection{Objetivos específicos}

%GUIA: as etapas que, somadas, conduzem ao objetivo geral; também
%iniciadas por verbo no infinitivo.

\begin{itemize}
\item enunciar cada etapa que, somada às demais, conduz ao objetivo
  geral;
\item iniciar cada item por verbo no infinitivo (levantar, implementar,
  testar\ldots);
\item manter entre três e cinco objetivos específicos, verificáveis ao
  final do trabalho.
\end{itemize}

\section{Estrutura do texto - Opcional}

Apresente aqui, caso já não o tenha feito no último parágrafo da
introdução, a organização do texto: um parágrafo indicando, capítulo a
capítulo, o que o leitor encontrará, por exemplo: ``o Capítulo~2
apresenta a fundamentação teórica; o Capítulo~3 discute os trabalhos
relacionados; o Capítulo~4 descreve a metodologia; o Capítulo~5
apresenta e discute os resultados; e o Capítulo~6 traz as conclusões e
os trabalhos futuros''. Utilize apenas um dos dois lugares (o parágrafo
final da introdução ou esta seção) e remova o outro.

\chapter{Fundamentação teórica}

%GUIA: embase, por meio de outros autores, os aspectos teóricos e
%conceituais da pesquisa. Apresente os principais autores sobre o tema
%e o tratamento que o assunto tem recebido da comunidade científica.
%Subdivida este capítulo em seções nomeadas de acordo com os assuntos
%tratados; não há limite mínimo ou máximo de seções — insira sempre
%texto entre uma seção e outra.
%
%REGRAS DE CITAÇÃO (NBR 10520:2023):
% - Citação INDIRETA (paráfrase das ideias do autor): sem aspas e sem
%   página obrigatória. Ex.: ... blablabla \cite{vonNeumann:1966:TSR}.
% - Citação DIRETA de até 3 linhas: no corpo do texto, entre "aspas
%   duplas", SEMPRE com indicação de página: \cite[p.~15]{chave}.
% - Citação DIRETA com MAIS de 3 linhas: use o ambiente quote (exemplo
%   abaixo), que já aplica recuo de 4 cm, fonte menor e espaço simples.
%   Não use aspas e informe a página.
% - CITAÇÃO DE CITAÇÃO (apud): quando não se teve acesso ao original.
%   Evite; priorize os textos originais. Use os comandos:
%     \apud{original}{consultada}  -> Moore (1979 apud Cormen et al., 1990)
%     \apudp{original}{consultada} -> (Moore, 1979 apud Cormen et al., 1990)
%   Na lista de Referências deve constar apenas a obra consultada.
% - O ponto final da frase vem DEPOIS do parêntese da citação:
%   correto:  ... blablabla (Silva, 2020).
% - Com 4 ou mais autores no .bib, a citação sai automaticamente como
%   "Silva et al." Use \cite{chave} para "(Silva, 2020)" e \citet{chave}
%   para "Silva (2020)" integrado à frase.
% - Toda fonte citada deve constar em bibliografia.bib (e vice-versa).

Apresente aqui os conceitos, as teorias e as tecnologias que sustentam
o trabalho, atribuindo sempre as ideias aos seus autores, como nesta
citação indireta, em que a fonte aparece ao final da
frase~\cite{Moore:1979:MAI}.

\section{Uma seção secundária}

Organize a fundamentação em seções temáticas, sempre com texto entre um
título e outro. Segundo \citet{vonNeumann:1966:TSR}, a citação
integrada à frase é feita com o comando \verb|\citet|; já a citação ao
final do período, entre parênteses, é feita com o comando
\verb|\cite|~\cite{Dijkstra:1968}.

\begin{quote}
  Use este ambiente para as citações diretas com mais de três linhas: o
  recuo de 4~cm, a fonte menor e o espaço simples exigidos pela norma já
  são aplicados automaticamente. Não use aspas, transcreva o texto
  exatamente como está na obra e informe sempre a página
  consultada~\cite[p.~10]{Nobre:2017}.
\end{quote}

Exemplo de citação de citação: conforme \apud{Moore:1979:MAI}{Cormen},
use o apud apenas quando não tiver acesso à obra original; nas
referências constará somente a obra efetivamente consultada
\apudp{Moore:1979:MAI}{Cormen}.

\subsection{Uma seção terciária}

Aprofunde aqui um tópico específico da seção anterior. Evite seções de
um único parágrafo e lembre-se: a norma exige texto entre o título de
uma seção e o título da seção seguinte.

\chapter{Trabalhos relacionados}

%GUIA: Apresente
%os sistemas, ferramentas ou estudos semelhantes ao seu (obtidos no
%levantamento bibliográfico), descrevendo cada um em uma seção ou
%parágrafo, SEMPRE com citação da fonte. Ao final, compare-os em um
%quadro (exemplo abaixo), incluindo o seu próprio trabalho na última
%linha: é essa comparação que evidencia a lacuna que o seu trabalho
%preenche e o seu diferencial. A atualidade e a relevância da
%bibliografia são critérios de avaliação da banca (Regulamento do TCC,
%art. 32). Evite usar somente comparações binárias (Sim - Não, Tem - Não Tem, etc).

Descreva aqui cada trabalho relacionado, um por parágrafo ou seção,
sempre com citação da fonte: o que o trabalho de \citet{Moore:1979:MAI}
faz, em que contexto foi aplicado e que resultados obteve.

Ao descrever a proposta de \citet{Cormen}, destaque o que a aproxima e
o que a diferencia do seu trabalho, bem como as limitações que o seu
trabalho pretende superar~\cite{vonNeumann:1966:TSR}.

O Quadro~\ref{quadro-trabrel} resume a comparação entre os trabalhos
relacionados e este trabalho.

\begin{quadro}[htbp]
\caption{Comparativo entre os trabalhos relacionados e este trabalho}
\label{quadro-trabrel}
\centering\begin{tabular}{|p{3.5cm}|p{3cm}|p{3cm}|p{4cm}|}
\hline
\textbf{Trabalho} & \textbf{Tecnologias} & \textbf{Público-alvo} & \textbf{Diferencial/limitação}\\
\hline
{\small \citet{Moore:1979:MAI}} & {\small Bla blabla} & {\small Bla} & {\small\em Bla blabla blablabla.}\\
\hline
{\small \citet{Cormen}} & {\small Bla blabla} & {\small Bla} & {\small\em Bla blabla blablabla.}\\
\hline
{\small \textbf{Este trabalho}} & {\small Bla blabla} & {\small Bla} & {\small\em Bla blabla blablabla.}\\
\hline
\end{tabular}
\begin{flushleft}
{\small Fonte: Elaboração própria.}
\end{flushleft}
\end{quadro}

\chapter{Metodologia}

%GUIA: apresente o conjunto de métodos e técnicas utilizados para a
%realização da pesquisa: a classificação do tipo de pesquisa, a
%caracterização do campo de análise (população e amostragem) e os
%instrumentos e procedimentos de coleta e de análise de dados.
%Em trabalhos de desenvolvimento de um sistema, descreva também: o método/processo de
%desenvolvimento adotado (ex.: iterativo, incremental, ágil), as
%tecnologias, linguagens, ferramentas e o ambiente utilizados e, quando
%houver experimentos, as métricas e o procedimento de avaliação.

Descreva aqui como a pesquisa foi realizada, com detalhamento
suficiente para que outro pesquisador possa reproduzi-la: a
classificação da pesquisa, o método ou processo de desenvolvimento
adotado, as tecnologias e ferramentas utilizadas e os procedimentos de
coleta e de análise de dados. Ilustre o processo quando isso ajudar o
leitor, conforme a Figura~\ref{figura}.

\begin{figure}[htbp]
\caption{Nome da figura}
\centering \includegraphics[scale=.4]{imagens/figura.pdf}\\
\begin{flushleft}
{\small Fonte: Elaboração própria.}
\end{flushleft}
\label{figura}
\end{figure}

%VARIANTE PARA TRABALHOS DE DESENVOLVIMENTO DE UM SISTEMA (Regulamento do
%TCC, art. 5º, natureza IV): descomente o capítulo abaixo. Ele apresenta
%a construção do sistema e antecede os Resultados.
%
%\chapter{Desenvolvimento}
%
%%GUIA: apresente a proposta e a construção do sistema: o levantamento
%%de requisitos funcionais e não funcionais (um quadro ajuda), a
%%modelagem (diagramas UML, modelo de dados), a arquitetura da solução
%%e os principais aspectos da implementação, ilustrados com figuras
%%(diagramas e telas do sistema). O código-fonte completo NÃO deve ser
%%transcrito na monografia: disponibilize-o em um repositório (ex.:
%%GitHub) e indique o endereço neste capítulo ou em apêndice.
%%Lembre-se: a entrega da cópia integral dos códigos-fonte é requisito
%%obrigatório de aprovação (Regulamento do TCC, art. 23) e, quando o
%%trabalho resultar em software, o registro no NIT deve ser encaminhado
%%(art. 23, §2º).
%
%Descreva aqui os requisitos funcionais e não funcionais, a modelagem,
%a arquitetura e os principais aspectos da implementação do sistema. O
%código-fonte completo do sistema está disponível em:
%\url{https://github.com/usuario/repositorio}.

\chapter{Resultados}

%GUIA: apresente os resultados da pesquisa. Em pesquisas com coleta de
%dados (naturezas I a III do art. 5º), analise os dados confrontando-os
%com os objetivos e as hipóteses/pressupostos, para confirmá-los ou
%rejeitá-los. Em trabalhos de desenvolvimento de um sistema (natureza IV),
%apresente os testes realizados (unitários, integração, aceitação), os
%cenários e resultados e a validação com usuários ou especialistas,
%quando houver, discutindo em que medida o sistema atende aos objetivos.
%Subdivida este capítulo em seções conforme a necessidade.

Apresente aqui os resultados obtidos e discuta-os à luz dos objetivos e
da pergunta/hipótese de pesquisa, pois resultados sem discussão não
demonstram que os objetivos foram alcançados. Use tabelas para dados
numéricos e quadros para comparações qualitativas, conforme a
Tabela~\ref{tabela} e o Quadro~\ref{quadro1}.

%REGRAS PARA ILUSTRAÇÕES, TABELAS E QUADROS (NBR 14724:2024):
% - Todo elemento deve ser citado no texto ANTES de aparecer (use \ref).
% - A indicação da FONTE, abaixo do elemento, é OBRIGATÓRIA, mesmo que
%   seja produção do próprio autor: use "Fonte: Elaboração própria." ou
%   "Fonte: Elaborado pelo próprio autor.". Se for de terceiros, cite a
%   fonte conforme a NBR 10520 (autor, ano) e referencie no final.
% - TABELA (dados numéricos/estatísticos): laterais abertas, use
%   \begin{table}. QUADRO (informações textuais/qualitativas): fechado
%   em todas as bordas, use \begin{quadro}.
\begin{table}[htbp]
\caption{Nome da Tabela}\label{tabela}
\centering\begin{tabular}{p{4cm}p{5cm}p{6cm}}
\hline
Blabla & Blabla & Blablabla\\
\hline
{\small Bla} & {\small Blabla} & {\small\em Bla blabla blablabla blabla
  blablabla blabla blablabla.}\\
{\small Bla} & {\small Blabla} & {\small\em Bla blabla blablabla blabla
  blablabla blabla blablabla.}\\
{\small Bla} & {\small Blabla} & {\small\em Bla blabla blablabla blabla
  blablabla blabla blablabla.}\\
\hline
\end{tabular}
\begin{flushleft}
{\small Fonte: Elaboração própria.}
\end{flushleft}
\end{table}

\begin{quadro}[htbp]
\caption{Nome do Quadro}\label{quadro1}
\centering\begin{tabular}{|p{4cm}|p{5cm}|}
\hline
Blabla & Blablabla\\
\hline
{\small Bla} & {\small\em Bla blabla blablabla blabla.}\\
\hline
{\small Bla} & {\small\em Bla blabla blablabla blabla.}\\
\hline
\end{tabular}
\begin{flushleft}
{\small Fonte: Elaboração própria.}
\end{flushleft}
\end{quadro}

\chapter{Conclusões finais e trabalhos futuros}

%GUIA: retome os objetivos e as hipóteses apresentadas na introdução,
%fornecendo evidências da solução do problema a partir dos resultados
%obtidos. Não inclua dados novos que não tenham sido apresentados
%anteriormente. Apresente, de maneira breve e ordenada, a síntese dos
%principais resultados e as contribuições trazidas pela pesquisa; no
%fechamento, pode-se apontar as limitações do método utilizado e
%recomendações ou sugestões para trabalhos futuros.

Retome aqui os objetivos e a pergunta/hipótese apresentados na
introdução e mostre, com base nos resultados, em que medida foram
alcançados ou respondidas. Sintetize as principais contribuições do
trabalho, reconheça as limitações do método e da solução e proponha, ao
final, trabalhos futuros que deem continuidade à pesquisa. Não apresente
neste capítulo dados ou informações que não tenham aparecido antes.

\section{Declaração sobre uso de Inteligência Artificial}

%GUIA (Portaria CNPq nº 2.664/2026 — Política de Integridade na
%Atividade Científica): a Portaria exige DECLARAR o uso de ferramentas
%de inteligência artificial generativa em QUALQUER fase da pesquisa
%(concepção, levantamento bibliográfico, escrita, análise de dados,
%submissão), especificando a ferramenta utilizada e a finalidade do uso.
%É VEDADO apresentar conteúdo gerado por IA como se fosse de autoria
%própria: o autor permanece integralmente responsável pelo conteúdo
%final, inclusive por plágio e por inexatidões produzidas pela
%ferramenta. Lembre-se ainda do Regulamento do TCC (art. 34): o uso de
%IA para a PRODUÇÃO do trabalho é considerado plágio e anula o TCC.

Declare aqui, de forma transparente, o uso de ferramentas de
inteligência artificial generativa em qualquer fase do trabalho,
conforme a Portaria CNPq nº 2.664/2026. Para cada uso, especifique:
(i)~a ferramenta, com nome e versão; (ii)~a finalidade do uso; e
(iii)~a fase do trabalho em que ocorreu. Caso nenhuma ferramenta tenha
sido utilizada, declare-o expressamente.

Exemplo de declaração: ``durante a elaboração deste trabalho, o autor
utilizou a ferramenta X (versão Y), com a finalidade de revisão
gramatical e de estilo, na fase de escrita. Após o uso, todo o conteúdo
foi revisado e editado pelo autor, que assume integral responsabilidade
pelo conteúdo final do trabalho''.

%EXEMPLOS DE TIPOS DE REFERÊNCIA (NBR 6023:2025) — o \nocite abaixo faz
%as entradas de demonstração aparecerem na lista de Referências SEM
%serem citadas no texto, apenas para você conferir a formatação de cada
%tipo. Tipos demonstrados (chave -> tipo):
%  Tanenbaum:2011     -> livro impresso
%  Goodfellow:2016    -> livro em meio eletrônico (gratuito no site oficial)
%  Pimentel:2011      -> capítulo de livro (coletânea com organizadores)
%  Dijkstra:1968      -> artigo de periódico com DOI como URL completa
%                        (citado na Fundamentação Teórica)
%  Michels:2024       -> trabalho publicado em anais de evento (SBIE)
%  Nobre:2017         -> trabalho de conclusão de curso/monografia
%                        (já citada no texto, na Fundamentação Teórica)
%  Costa:2019         -> dissertação (mestrado)
%  Garcia:2014        -> tese (doutorado, em repositório institucional)
%OBS.: TODAS as entradas são obras reais, verificadas nas fontes.
%  ABNT:25010:2011    -> norma técnica
%  Brasil:LGPD:2018   -> legislação em meio eletrônico
%  CNPq:2026          -> portaria de órgão público
%  Python:2025        -> software
%  StackOverflow:2025 -> página de site na internet
%  BernersLee:2009    -> vídeo online (palestra TED)
%  Guanabara:2024     -> podcast, entrevistado como autor (NBR 6023:2025)
%NO SEU TRABALHO: remova esta linha \nocite e as entradas não usadas do
%bibliografia.bib — a lista deve conter somente obras citadas no texto.
\nocite{Tanenbaum:2011,Goodfellow:2016,Pimentel:2011,Michels:2024,Costa:2019,Garcia:2014,ABNT:25010:2011,Brasil:LGPD:2018,CNPq:2026,Python:2025,StackOverflow:2025,BernersLee:2009,Guanabara:2024}

\bibliographystyle{abnt}
\bibliography{bibliografia}

% Bibliografia http://liinwww.ira.uka.de/bibliography/index.html um
% site que cataloga no formato bibtex a bibliografia em computacao

%REGRAS PARA AS REFERÊNCIAS (NBR 6023:2018) — preencha bibliografia.bib:
% - Separe os autores com " and " no campo author e escreva cada um no
%   formato "Último sobrenome, Inicial do primeiro nome". Exemplo:
%   Alex Mulattieri Suarez Orozco -> author = "Orozco, A."
%   O estilo formata sozinho: OROZCO, A. (e separa por ponto e vírgula).
% - Entrada da referência por tipo de sobrenome:
%     simples: entrada pelo último ..................... SILVA, A.
%     composto (substantivo+adjetivo): mantém os dois
%       (no .bib: "{Castelo Branco}, Ana") ............. CASTELO BRANCO, A.
%     prefixo de/da/do/von/van: vai APÓS o nome, não
%       integra o sobrenome (no .bib: "de Mota, Ana") .. MOTA, A. de
%     Filho/Neto/Júnior: integra o sobrenome
%       (no .bib: "{Oliveira Filho}, João") ............ OLIVEIRA FILHO, J.
%     com hífen: sobrenome único ....................... TELES-PONTES, G.
%     feminino com sobrenome do cônjuge: manter a forma
%       usada habitualmente pela autora
%     estrangeiro: convenção do país de origem
%       (no .bib: "{García López}, María") ............. GARCÍA LÓPEZ, M.
%     chinês/japonês/coreano: sobrenome já vem primeiro
%       (no .bib: "Mao, Zedong") ....................... MAO, Z.
% - Para documentos online, preencha url e urldate no formato
%   AAAA-MM-DD (ex.: urldate = {2026-07-06}); o "Disponível em:" e o
%   "Acesso em:" são gerados automaticamente no formato da norma.
% - Só liste obras efetivamente citadas no texto.
% - Padronize nome do autor, ou SOBRENOME, Iniciais; ou SOBRENOME, Nome;
%   Não misture as duas formas de representar o autor.



% Apêndices (Opcional) - Material produzido pelo autor
%REGRA (NBR 14724:2024): fontes citadas EXCLUSIVAMENTE em apêndice ou
%anexo devem ser referenciadas no próprio apêndice/anexo, e NÃO na lista
%de referências principal. Este é um dos itens que mais gera devolução
%do trabalho pelas bibliotecas.
%COMO FAZER: envolva o conteúdo do apêndice/anexo no ambiente
%referenciasapendice (exemplo abaixo). As citações \cite/\citet funcionam
%normalmente e a lista de referências dessas obras é gerada
%automaticamente ao final do próprio apêndice.
\apendices
\chapter{Título do apêndice}

\begin{referenciasapendice}
Os apêndices apresentam material de apoio elaborado pelo próprio autor,
como questionários, termos de consentimento, diagramas completos ou
manuais. Cite as fontes normalmente~\cite{Knuth:1984}: as obras citadas
somente no apêndice são listadas logo abaixo, no próprio apêndice.
\end{referenciasapendice}

% Anexos (Opcional) - Material produzido por outro
\anexos
%GUIA: os ANEXOS são documentos NÃO elaborados pelo autor, que servem de
%fundamentação, comprovação e/ou ilustração (já os APÊNDICES são de
%autoria própria). Ambos são identificados por letras maiúsculas
%consecutivas, travessão e o respectivo título.
\chapter{Título do anexo}

Os anexos apresentam material de apoio \textbf{não} elaborado pelo
autor, utilizado como fundamentação, comprovação ou ilustração, por
exemplo, normas, leis, manuais de terceiros ou documentos da organização
estudada.

Identifique cada anexo por letra maiúscula consecutiva, travessão e
título, e cite no corpo do texto a origem do material aqui reproduzido.

\chapter{Título do outro anexo}

Utilize um novo anexo para cada documento de origem distinta.

Remova deste modelo os apêndices e anexos que não forem utilizados no
seu trabalho.

\end{document}

